\begin{table}[H]
    \centering
    \caption{Baseline Estimates under Treatment-Structure Clustering}
    \label{tab:treatment_structure_clustering}
    \resizebox{0.98\textwidth}{!}{%
    \begin{tabular}{lcccc}
        \toprule
        Clustering & \(IMS\) & \(IMS\_D\) & Clusters & N \\
        \midrule
        (1) Firm & -0.4829** (0.2362) & -0.1434* (0.0794) & 4604 & 15,381 \\
        (2) Audit firm & -0.4829** (0.2428) & -0.1434 (0.0958) & 96 & 15,381 \\
        (3) Audit-firm–industry & -0.4829** (0.2277) & -0.1434* (0.0839) & 1205 & 15,381 \\
        (4) Firm + audit firm & -0.4829* (0.2470) & -0.1434 (0.0970) & 4604 × 96 & 15,381 \\
        (5) Audit firm + industry & -0.4829* (0.2770) & -0.1434 (0.1011) & 96 × 47 & 15,381 \\
        (6) Industry & -0.4829* (0.2639) & -0.1434 (0.0900) & 47 & 15,381 \\
        \bottomrule
    \end{tabular}%
    }
    \par\smallskip
    \begin{minipage}{0.98\textwidth}
        \footnotesize
        \textit{Notes:} Same baseline specification with year and industry fixed effects. Only the standard-error clustering changes. Two-way rows report Cameron--Gelbach--Miller multiway clustering. Coefficient estimates are followed by standard errors in parentheses. ***, **, and * denote two-tailed significance at 1, 5, and 10 percent.
    \end{minipage}
\end{table}
